\chapter{Live Service Upgrade}

CharlotteOS can replace a running service without losing its state. This is
\textbf{restart-with-state}, not true zero-downtime: in-flight requests to the
old instance fail during the handoff window and clients must retry. But the
new instance resumes with the old instance's state intact, and clients
reconnect transparently through the name service's generation tracking.

\section{The four primitives}

Live upgrade builds on four primitives already implemented in the kernel and
name service:

\begin{enumerate}
    \item \textbf{Memory-object ownership transfer.} The kernel's
        \texttt{memory::object::move\_to} moves a page-backed memory object
        from one address space to another. The sender loses access; the
        receiver gains it. This is the vehicle for service state.

    \item \textbf{Generation tracking.} The userspace name service records an
        instance generation that increments on restart. A client calling
        through a stale connection from the previous generation receives
        \texttt{EndpointClosed}. Re-looking up the name returns the new
        generation and a fresh connection to the new instance.

    \item \textbf{EndpointClose and Cancelled.} When a service domain exits,
        the kernel closes its endpoint. Queued requests that were never
        delivered complete as \texttt{EndpointClosed}. Delivered calls whose
        server exited before replying complete as \texttt{Cancelled}. Clients
        see a deterministic failure, not a hung request.

    \item \textbf{Bootstrap capability delivery.} The supervisor writes
        initial capabilities to the config page before a domain starts. For a
        live upgrade, the bootstrap contract is extended: in addition to the
        name service connection, the new domain receives the old service's
        state (as memory objects) and the old service's endpoint capability.
\end{enumerate}

\section{The handoff protocol}

\begin{verbatim}
     Old Service         Supervisor         New Service
          |                    |                   |
     OP_HANDOFF               |                   |
          |-------------------->                   |
          |   drain in-flight, |                   |
          |   serialize state  |                   |
          |   -> move_to       |                   |
          |<-------------------|                   |
          |  reply: "ready"    |                   |
          |                    |  spawn_upgrade()  |
          |  (exit)            |------------------>|
          |                    |  bootstrap:       |
          |                    |    ns_connection  |
          |                    |    state_memory   |
          |                    |    endpoint_cap   |
          |                    |                   |
          |                    |         new service reads state,
          |                    |         takes over endpoint,
          |                    |         registers (new generation)
          |                    |                   |
                    clients see EndpointClosed
                    -> re-lookup -> fresh connection
                    -> resume from where they left off
\end{verbatim}

The old service receives an \texttt{OP\_HANDOFF} request (typically from the
supervisor or a service manager). It drains any in-flight work, serialises its
mutable state into memory objects, moves those objects to the supervisor's
address space, and replies ``ready'' before exiting.

The supervisor then spawns the new service instance with an extended bootstrap
contract: the name service connection, the handoff state memory objects, and
the old service's endpoint capability. The new service deserialises the state,
takes over the endpoint, and re-registers under the same name. The name
service bumps the generation, invalidating all stale client connections.

\section{What the supervisor provides}

\begin{verbatim}
pub struct UpgradeGrant {
    pub state_caps: Vec<MemoryObjectCap>,
    pub endpoint_cap: CapabilityId,
}

pub fn spawn_upgrade(
    image: &[u8],
    name_service: &NameServiceHandle,
    old_domain: ServiceDomain,
    grant: UpgradeGrant,
) -> ServiceDomain;
\end{verbatim}

\texttt{spawn\_upgrade} loads the new service ELF, moves the state memory
objects to the new domain, delegates the old endpoint cap to the new domain,
writes the extended bootstrap config page, starts the new domain, and tears
down the old domain (reclaiming all remaining resources).

\section{What the service must implement}

The old service must handle \texttt{OP\_HANDOFF}:

\begin{verbatim}
OP_HANDOFF => {
    // Stop accepting new work.
    // Serialize all mutable state into memory objects.
    // Move those objects to the supervisor's address space
    //   (the supervisor passes its ASID in the handoff call).
    // Reply "ready", then exit.
}
\end{verbatim}

The new service reads the extended bootstrap contract:

\begin{verbatim}
fn main(ctx: Context) -> ! {
    let ns_connection = ctx.bootstrap_cap();
    let state_count = ctx.handoff_count();
    let state = ctx.handoff_state_cap();
    let endpoint = ctx.handoff_endpoint_cap();

    // Deserialize state, register the endpoint under the
    // same name (NS bumps generation), start serving.
}
\end{verbatim}

\section{What clients see}

A client holding a stale connection from the previous generation sees a clean
failure and retries:

\begin{verbatim}
// Before upgrade:
let conn = ns.lookup("payroll")?;       // generation 3
payroll.call(OP_CALCULATE, employee_id); // in-flight

// During upgrade: call completes as EndpointClosed.
// Client retries:
let conn = ns.lookup("payroll")?;       // generation 4
payroll.call(OP_CALCULATE, employee_id); // reaches new instance
\end{verbatim}

The new instance has the old instance's state, so it can resume processing
where the old instance left off. From the client's perspective, a single
request failed and was retried --- the rest of the session continues
uninterrupted.

\section{Restart-with-state vs.\ zero-downtime}

This design is \textbf{restart-with-state}: in-flight requests to the old
instance fail during the handoff window, and clients must retry. This is
sufficient for the vast majority of services --- the handoff window is
measured in milliseconds, and clients already handle transient failures
through the name service's generation tracking.

True zero-downtime would require one of:

\begin{itemize}
    \item \textbf{Queue migration:} the old service's endpoint queue would
        need to be transferred to the new service atomically, so no request
        is ever rejected. This is a future kernel feature.
    \item \textbf{Sidecar model:} the new service starts alongside the old
        one, both sharing the same endpoint, and the old service stops
        accepting new work. The kernel does not currently support multiple
        owners for one endpoint.
\end{itemize}

\section{Relationship to crash recovery}

The live upgrade path is the graceful variant of the crash-recovery path
already exercised by the UART driver test (Chapter~12). In both cases, the
supervisor reclaims device authority, invalidates stale connections, and
spawns a new instance with a bumped generation. The difference is that live
upgrade transfers state explicitly through memory objects before the old
instance exits, rather than starting the new instance from scratch.

\section{Why this matters for critical software}

\begin{enumerate}
    \item \textbf{Services can be upgraded without losing state.} A bug fix
        or feature update to a filesystem, network stack, or consensus node
        does not require a system reboot. The new binary starts with the old
        binary's state.

    \item \textbf{The upgrade window is bounded.} The old service drains
        in-flight work, serialises state, and exits. The new service starts
        and registers. Clients retry. The entire sequence completes in
        milliseconds.

    \item \textbf{Failure during upgrade is safe.} If the new service fails to
        start, the supervisor can restart the old service (or a fallback).
        The state memory objects are owned by the supervisor during the
        transition, so they are not lost.

    \item \textbf{The building blocks are already proven.} Memory-object
        transfer, generation tracking, deterministic teardown, and bootstrap
        delivery are all exercised by the existing self-test suite. The live
        upgrade protocol adds approximately 150 lines of orchestration code.
\end{enumerate}

\endinput
